%% The technical contributions, at the end of the introduction, input by
%% paper.tex. One sentence each (two in a pinch); each is novel (to our knowledge, per
%% notes/novelty-literature.md: C2-C4, C6, C7, C9) and is expanded in the
%% section it points to. The section labels are placeholders until those
%% sections exist.

Our technical contributions are:
\begin{itemize}
\item A proof, mechanized in Agda, that a ten-pass compiler from a
  call-by-value $\lambda$-calculus with iso-recursive types, through
  closure conversion, to a typed assembly language is correct, using a
  graph-model semantics for every intermediate language. Most passes are
  proved by an equation between the denotations of a program before and
  after the pass (\S\ref{sec:compiler-proof}).
\item Approximant logical relations, which relate two denotational
  semantics by recursion on the depth of a finite approximation of one
  of them, and with which we prove the reflection direction of lowering
  closures into code and environment pairs, where recursive types rule
  out recursion on types (\S\ref{sec:approx-lr}).
\item Proofs that the graph-model semantics of the source language and of
  the assembly language are adequate with respect to interpreters for
  those languages, which yield an end-to-end theorem that the source
  program and the compiled program return the same integer or both
  diverge (\S\ref{sec:adequacy}).
\item A comparison with proofs of the same compiler in three operational
  styles (abstract machines, clocked interpreters and reduction
  semantics), which shows that the denotational proofs are among the
  smallest when program values stay at the same level of abstraction and
  are the largest when values are lowered (\S\ref{sec:evaluation}).
\end{itemize}
